\documentclass[../main.tex]{subfiles}
\begin{document}
%==========================================TIÊU ĐỀ TẬP========================================
\begin{table}[h]
    \begin{adjustwidth}{-1cm}{}
        \begin{tabular}{>{\centering\arraybackslash}p{6cm}|>{\raggedright\arraybackslash}p{12cm}}
            \multirow{5}{6cm}
            % Cột bên trái: ÉPISODE và số tập
            {\\[-40pt]
            \centering
                {\fontsize{20pt}{20pt}\selectfont{\outlinetext[1]{vol}{white}{\flaregothic PERSONNALISÉ}}\\[12pt]
                {\fontsize{36pt}{36pt}\selectfont{\outlinetext[1]{vol}{white}{\vnmdisp SÁU}}}\\[6pt]     % use for custom episodes
                {\fontsize{20pt}{20pt}\selectfont{\outlinetext[1]{vol}{white}{\cafeteria (C-6)}}}\color{black}}
                \\[18pt]
                % {\fontsize{14pt}{14pt}\selectfont{\outlinetext[1]{vol}{white}{\flaregothic [I]}}}\\
                % {\fontsize{18pt}{18pt}\selectfont{\outlinetext[1]{vol}{white}{\flaregothic [II]}}}\\
                \color{black}
            }
             &
            % Dòng thứ nhất: tên tập tiếng Nhật
            {\fotskip\fontsize{18pt}{18pt}\selectfont そろそろ\ruby{一人}{ひとり}で\ruby{抱}{かか}え\ruby{込}{こ}むの、やめてみませんか？}                                                                                                                                                                         \\[6pt]
            % Dòng thứ hai: tên tập tiếng Anh
             & {\cafeteria\fontsize{18pt}{18pt}\selectfont A Shift in the New Normal}                                                                                                                                                                               \\[4pt]
            % Dòng thứ ba: tên tập tiếng Việt
             & {\vnmsans\fontsize{18pt}{18pt}\selectfont Bản báo cáo định kỳ}                                                                                                                                                                         \\[8pt]
            % Dòng thứ tư: Nhân vật xuất hiện
             & {\caviar{\fontsize{14pt}{14pt}\selectfont Nhân vật: \emp{kuon2} \emp{ryuuhou2} \emp{kaigou2} (hồi tưởng) \emp{game} \emp{achan}}} \\[8pt]
            % Dòng cuối cùng: Thời gian diễn ra của tập (thường sẽ là ngày phát hành)
             & {\continuum\fontsize{16pt}{16pt}\selectfont Ngày phát hành: 03/12/2021}                                     \\[18pt]
        \end{tabular}\\[8pt]
    \end{adjustwidth}
\end{table}
%=========================================TRUYỆN CHÍNH========================================
\color{black}

\txtbx[2]
{{\color{vol} \uline{Tóm tắt:}} Kuon báo cáo định kỳ với A-chan về hiệu suất làm việc và thiệt hại do holomem gây ra - một danh mục luôn khiến cậu đau đầu. A-chan đề nghị Kuon chuyển sang vị trí đào tạo nhân sự mới để giảm gánh nặng. Kuon ban đầu từ chối vì sợ mất lương, nhưng sau những trải nghiệm qua các vũ trụ, cậu nhận ra gia đình và hạnh phúc mới là quan trọng nhất, nên đồng ý. Ryuuhou cũng nhận lời thử việc tại COVER, đảm nhiệm mảng luyện thanh và biên đạo. Cả hai cùng nhau thảo luận kế hoạch ``uốn nắn'' Kaigou cho vị trí quản lý \hls. Kết thúc buổi họp, Kuon và Ryuuhou bước ra phố, Ryuuhou tiết lộ lý do cô làm việc với cậu - cốt là để thể hiện sự trưởng thành và mong muốn sánh vai cùng anh trai.}

\pov{Hikawa Kuon}

Thấm thoát đã khoảng một tuần trôi qua kể từ khi những giông bão và biến cố dồn dập từ sự kiện đa vũ trụ kia chính thức khép lại. Mọi xáo trộn dần lắng xuống, nhường chỗ cho cái nắng sớm dịu dàng của một ngày đầu tháng cuối cùng của năm 2021. Căn nhà nhỏ của gia đình Hikawa rốt cuộc cũng tìm lại được bầu không khí ấm áp vốn có, hay chính xác hơn là một ``cuộc sống bình thường mới'', nơi nhịp sinh hoạt thường nhật của cả gia đình tôi - cả những người cũ lẫn những người mới - bắt đầu đi vào quỹ đạo ổn định.

Chuyến hành trình của tôi vốn rất dài dằng dặc và nhiều chi tiết, thế nhưng tôi có thể gói gọn lại trong vài điểm chính như sau:

Mọi chuyện bắt nguồn từ văn phòng \hll, khi một trục trặc ma thuật phát sinh từ những bản sao của Shion đã vô tình tạo ra một lỗ xoáy không gian, khiến tôi bị cuốn vào đó. Tôi tỉnh lại giữa The Void, một khoảng không vô tận nằm ngoài dòng chảy thời gian, tại đây tôi đã có cơ hội gặp gỡ và giúp hai chị em Kaigou cùng Suzuno đoàn tụ với nhau. Ngay sau đó, cả ba chúng tôi cùng rơi vào thế giới số của ``There Is No Game'' và kết thân với Game, một trí tuệ nhân tạo đơn độc đã quyết định cùng chúng tôi vượt qua các không gian thứ nguyên khác nhau.

Tập thể bốn người chúng tôi tiếp tục trải qua mười hai ngày đầy biến động nhưng vô cùng ấm cúng tại thế giới Nijisanji, trước khi một lỗ xoáy mới lại cuốn chúng tôi đến thị trấn mịt mù sương Aogami. Tại đây, trong thế giới của \hll\ ERROR, toàn bộ ký ức của chúng tôi bị xóa sạch, buộc phải mang danh tính giả để đối mặt với vòng lặp bi thương của Misora Shino - tức là Sora của thế giới này - cùng với lời nguyền hiến tế đã tồn tại từ thuở xa xưa. Nhờ vào sự kiên định không bỏ cuộc, chúng tôi đã cử hành nghi thức thanh tẩy để giải thoát linh hồn của Naomi, giúp cô bé đã bước sang tuổi thứ 1300 này được tái sinh thành một con người trọn vẹn bằng xương bằng thịt.

Khi tất cả những ``lỗi lầm'' kia đều được ``sửa chữa'', một lỗ xoáy trắng đen đưa cả nhóm trở về văn phòng \hll\ vào đúng thời khắc tôi vừa rời đi ban đầu. Sau khi về nhà, một kết quả xét nghiệm ADN gây sốc đã xác nhận Kaigou, Suzuno và Naomi đều có quan hệ huyết thống với gia đình Hikawa, là hệ quả của sự cộng hưởng giữa các thứ nguyên. Ba cô gái chính thức được ghi tên vào sổ hộ khẩu, biến lời hứa về một ``mái nhà thực sự'' mà tôi đã nói tại The Void thành hiện thực về mặt pháp lý.

\model{25}{l}{5cm}{kuon2}

Những ngày tiếp theo, Kaigou trở thành quản lý tại \textit{\hls} - một nhánh các VTuber chỉ gồm toàn nam giới, trong khi Suzuno và Naomi bắt đầu cuộc sống học đường bình thường như bao bạn bè đồng trang lứa. Thực thể Game vẫn tiếp tục tồn tại bên trong chiếc đồng hồ đeo tay của tôi, trở thành một thành viên không thể thiếu trong đại gia đình chúng tôi. Cuộc phiêu lưu xuyên vũ trụ này đã thay đổi hoàn toàn cuộc đời tôi, biến tôi từ một quản lý đơn độc thành người anh cả của bốn cô em gái trong một mái nhà luôn tràn ngập tiếng cười và hạnh phúc.

Đó là toàn bộ những điểm chính tôi muốn tóm tắt. Nếu các bạn muốn tìm hiểu kỹ hơn về từng giai đoạn, hãy tìm đọc lại các quyển **V: Những Mảnh Vỡ Của Thời Gian**, **NJ1: Mười Hai Ngày ở Thế Giới Cầu Vồng** và **ER: Lời Nguyền ở Thị Trấn Aogami**.

Để đánh dấu cho một khởi đầu mới mẻ và nhiều năng lượng hơn, sáng hôm nay, cả tôi và Ryuuhou đều quyết định chọn cho mình những bộ trang phục mới trước khi bước ra khỏi nhà. Tôi khoác lên mình chiếc áo sơ mi kẻ sọc đỏ đen khoẻ khoắn phối cùng quần đen gọn gàng, thắt dây lưng nâu và tất nhiên không thể thiếu đôi xăng-đan quen thuộc đã đi cùng tôi qua bao nhiêu mùa mưa nắng. Đi cùng với đó là chiếc đồng hồ đeo tay điện tử quen thuộc đã đồng hành cùng tôi suốt hơn hai năm qua, và cũng là nơi mà trợ lý ảo của tôi - Game - trú ngụ.

Thế nhưng, điều khiến tôi phải sững người mất vài giây chính là sự xuất hiện của Ryuuhou.

Con bé bước ra từ phòng ngủ trong một diện mạo hoàn toàn tươi mới: chiếc áo thun màu cam nổi bật có in hình đám mây nhỏ trước ngực, kết hợp cùng chiếc váy denim năng động dài tới giữa đùi, đi cùng với đôi tất trắng cao đến đầu gối và đôi giày thể thao trắng tinh tươm. Chiếc cài tóc hình chữ X màu cam vẫn được cài lệch ở trên mái đúng vị trí như cũ. Trông bề ngoài con bé không khác nào một cô nàng tuổi teen đúng nghĩa.

\model{25}{r}{9.5cm}{ryuuhou2}

Điều đáng ngạc nhiên ở đây không chỉ là việc con bé ăn mặc có phần khá phong phanh giữa cái thời tiết dạo này đã bắt đầu trở lạnh đầu tháng Mười Hai, mà chính là chiếc váy denim kia. Bình thường, con bé cực kỳ hiếm khi chịu mặc váy. Cả năm dài tháng rộng, em ấy chỉ trung thành tuyệt đối với quần đùi hoặc quần dài tùy theo mùa. Lần gần đây nhất mà tôi nhớ Ryuuhou chịu mặc đầm là trong đám cưới của một người họ hàng bên nội từ cách đây khoảng\dots\ tôi không biết nữa, ba tháng trước?

Và cứ mỗi lần tôi hỏi, con bé đều chun mũi viện lý do láu cá rằng: ``Mặc váy làm gì để ông anh trai nhìn trộm à?''. Nhưng thực chất, tôi biết tỏng đó chỉ là cách con bé che đậy cái tính cách tomboy nghịch ngợm, thích bay nhảy chốn đông người của mình mà thôi.

Nhìn cô em gái cứ loay hoay nắm lấy gấu váy nhấp nhổm, tôi không khỏi cất tiếng hỏi: ``Ryuuhou này, dạo này thời tiết bắt đầu trở lạnh rồi đấy. Mọi khi anh thấy em có bao giờ chịu mặc váy đâu, toàn chui vào quần dài cho thoải mái. Sao hôm nay tự nhiên lại `đổi gió' thế này?''

Ryuuhou khẽ giật mình, hai gò má hơi hẩy hồng lên. Con bé lấy tay gãi gãi mái tóc ngắn, ngượng ngùng hất cằm đáp: ``Thì\dots\ thì gia đình mình vừa trải qua bao nhiêu chuyện sóng gió, giờ mọi thứ ổn định lại rồi nên em cũng muốn có một `khởi đầu mới' chứ sao! Với lại hôm nay em đi theo anh lên công ty mà, phải ăn mặc chỉn chu và nữ tính một chút cho anh mát mặt chứ. Anh có biết em đã phải gom hết dũng khí tích góp cả tuần nay mới dám chui vào cái váy này không hả?''

Tôi không nghĩ cái ``khởi đầu mới kia'' là lý do duy nhất để em tôi dám phá cách như thế này. Hẳn phải có lý do nào đó khác mà tôi cũng chẳng tài nào biết được.

Vừa khi con bé dứt câu, chiếc đồng hồ đeo trên cổ tay tôi bèn lóe lên một tia sáng xanh nhè nhẹ, kèm theo cái giọng nói trầm đục, vốn luôn hay trêu chọc người khác đã quá quen thuộc, xen ngang: ``Tôi nghĩ cậu đừng nên phí sức miệt mài phân tích tâm lý các cô gái ở lứa tuổi mới lớn, dù hệ thống của tôi có chạy đến mười nghìn vòng thuật toán cũng chẳng tài nào giải mã nổi đâu.''

Ngay lập tức Ryuuhou tròn xoe mắt, người cúi gập xuống, ngón tay chỉ thẳng vào mặt đồng hồ hét lên: ``Game-san! Lại còn cả ông đi dòm ngó và nói móc tôi nữa rồi đấy!''

``Tôi chỉ phân tích số liệu thực tế mà cảm biến thu thập được thôi, Ryuuhou-user,'' ông ấy vẫn thản nhiên trả lời qua chiếc loa nhỏ. ``Nhắn cũng phải ghi nhận, lòng dũng cảm của cô hôm nay thật đáng ngưỡng mộ, ít nhất thì cho đến khi cơn gió lạnh mùa đông kéo tới, khiến cô phải run cầm cập vì đóng băng ngoài đường.''

Ryuuhou lập tức bật lại, giọng đầy bực bội: ``Cái ông này hay thật! Tôi có ốm đau hay thế nào thì có phải đến lượt ông lo đâu, cứ ở lì trong đồng hồ mà làm nhiệm vụ của mình đi! Đừng có mà ngày ngày chỉ biết đi bới móc chuyện của người khác!''

\dots Thôi thì, nếu con bé chẳng thấy phiền muộn gì với chuyện ăn mặc thế kia thì tôi cũng chẳng có lý do gì để phản đối. Tôi chỉ thầm cầu mong rằng đừng có lúc nào con bé đột nhiên la hét lên, gọi tôi là kẻ biến thái chỉ vì tôi vô tình trông thấy\dots\ ờm, mấy thứ bên trong váy của con bé.

Tôi chỉ bật cười lắc đầu trước màn khẩu chiến quen thuộc giữa ông chú già đời này và cô em gái nhỏ. Sự náo nhiệt và năng lượng tích cực ấy suốt quãng đường dẫu sao cũng giúp tôi vơi bớt phần nào áp lực đang chực chờ phía trước. Bởi vì đối với một người gánh vác khối lượng công việc vận hành đồ sộ như tôi, ngày đầu tiên của mỗi tháng chẳng khác nào một ngày ``phán quyết''.

Cứ đến hẹn lại lên, nhiệm vụ quen thuộc của tôi là ngồi tổng hợp lại toàn bộ khối lượng công việc khổng lồ đã hoàn thành trong tháng qua: từ lịch trình hoạt động, tiến độ các dự án âm nhạc, cho đến công tác hậu cần của toàn bộ các thành viên. Thế nhưng, điều đó vẫn chưa phải là phần tàn khốc nhất.

Thứ chiếm nhiều tâm trí, thời gian và khiến tôi phải bóp trán đau đầu nhất trong bản báo cáo sắp sửa nộp cho A-san chính là\dots\ danh mục thống kê thiệt hại do các các ``nàng hề'' gây ra cho văn phòng.

Từ những vụ lỡ tay làm hỏng thiết bị phòng thu đắt tiền, đập phá đạo cụ livestream, cho đến mấy sự cố hy hữu bộc phát năng lượng khiến linh kiện điện tử cháy đen hay bàn ghế bị gãy đôi\dots\ tất cả đều được tôi tỉ mỉ ghi chép không sót một xu. Mỗi dòng con số thiệt hại tăng lên trong bảng thống kê ấy là một lần tôi cảm thấy ruột gan mình đau như cắt thay cho ngân sách công ty, đồng thời tự hỏi rốt cuộc mình đang làm một quản lý tài năng hay đang làm thủ kho dọn dẹp hậu quả cho một tập thể toàn những ``tai họa di động''.

Tôi ôm chặt xấp tài liệu nặng trịch trong tay, hít một hơi sâu khi đứng trước cửa văn phòng COVER. Báo cáo tháng này hứa hẹn sẽ lại là một buổi làm việc đầy dở khóc dở cười.

\ct

Cánh cửa phòng làm việc của A-san bật mở với một tiếng cạch nhẹ. Cô nàng Đeo kính - vừa là đồng nghiệp, vừa là đàn chị của tôi - đang lọt thỏm giữa hàng đống hồ sơ, vừa thấy chúng tôi bước vào liền ngẩng đầu lên, nhẹ nhàng đẩy gọng kính rồi nở nụ cười chào đón. Ryuuhou nhanh nhẩu kéo ghế ngồi xuống bên cạnh, còn tôi thì cẩn thận đặt xấp báo cáo dày cộp lên mặt bàn.

A-san lật qua từng trang tài liệu, ánh mắt lướt nhanh qua những dòng số liệu chi tiết đến từng yên. Cô gật gù đánh giá: ``Ấn tượng thật đấy, Kuon-san. Hiệu suất làm việc tháng này của em vẫn giữ ở mức tuyệt đối như mọi khi. Mọi sổ sách, chi phí phát sinh cho đến lịch trình của các nhánh đều được sắp xếp không chê vào đâu được.''

Nhưng rồi, nụ cười trên môi cô chầm chậm hạ xuống. Chị tựa lưng vào ghế, thở dài một hơi đầy ái ngại khi nhìn gạt sang xấp lịch làm việc dày đặc của tôi: ``Thế nhưng mà\dots\ nhìn vào đây chị thật sự thấy lo cho em đấy. Năm nay em 25 tuổi rồi, gắn bó với công ty cũng đã hơn hai năm rưỡi. Khối lượng công việc vận hành ở công ty vốn đã ngập đầu, vậy mà bài vở, đồ án trên trường của em gần này cũng nặng không kém. Rốt cuộc em xoay sở kiểu gì để không gục ngã hay vậy?''

Tôi gãi gãi gáy, nở nụ cười xoa dịu: ``Thì em cứ phân chia thời gian hợp lý thôi chị. Ban ngày lên lớp, tối về xử lý việc công ty, rảnh lúc nào thì chợp mắt lúc đó. Em vẫn còn sức khỏe tốt mà, chị không phải lo quá đâu.''

A-san lắc đầu, khóe miệng trồi lên nửa nụ cười bất lực khi ánh mắt xuyên qua tròng kính dán chặt vào tôi: ``Em cứ luôn quay đi quay lại một câu `vẫn khỏe, vẫn ổn' được mãi à. Cả sếp Tanigo lẫn con bé Ryuuhou bên cạnh đây đều đã than thở với chị không biết bao nhiêu lần về cái thói ôm đồm việc, ngó lơ sức khỏe của em đấy.''

Ngay lập tức, cô em gái liền nhăn mặt bật vào, đạp nhẹ vào chân tôi dưới gầm bàn như thể để chứng minh lời của A-san: ``Đúng đó chị A! Hôm qua em còn dậy lúc ba giờ sáng mà thấy Onii vẫn ngồi trước máy tính gõ lạch cạch, gọi đi ngủ thì cứ ỉ ôi bảo `sắp xong rồi', thế mà cho đến khi em đi học vẫn thấy anh ấy cắm mặt vào màn hình!''

Chưa hết lời, chiếc đồng hồ trên cổ tay tôi cũng bật sáng, giọng Game thản nhiên vang lên xen vào: ``Dữ liệu lịch sử ghi nhận Kuon-user có những sáu ngày trong tuần vừa qua ngủ ít hơn năm tiếng. Một tên cuồng công việc thực thụ, đúng như các hệ thống tôi đã phân tích từ lâu.''

Nói đến đây, A-san dừng lại một chút, ánh mắt trở nên trầm ngâm hơn. Cô chống tay lên bàn, nhìn thẳng vào mắt tôi: ``Chị hỏi thật nhé Kuon\dots\ Tất cả những gì em đang ra sức gánh vác, cố làm quần quật không nghỉ tay như thế này\dots\ rốt cuộc cũng chỉ vì một chữ `tiền' thôi đúng không?''

Không khí trong phòng chìm vào khoảng lặng ngắn. Chị ấy tiếp lời, giọng điệu đầy thắc mắc: ``Theo đúng thỏa thuận ban đầu, số tiền lương công ty trả cho em vốn rất đều đặn, thậm chí tiền thưởng làm thêm giờ còn nhiều tới mức quá dư dả để em đóng học phí cho cả kỳ. Chỉ cần làm chưa đầy hai tháng là em đã dư sức chi trả cho một kỳ học kéo dài tận ba tháng rồi. Thế mà sao lúc nào chị cũng cảm thấy em trong trạng thái thiếu thốn và áp lực tiền bạc đến thế?''

Tôi thu lại nụ cười xã giao, chậm rãi khoanh tay trước ngực. Nhìn ra ngoài cửa sổ kính của văn phòng, tôi đáp lại bằng một giọng điệu vô cùng gọn gẽ và trầm lắng: ``Chị sẽ không biết những gì sẽ xảy đến trước mắt đâu.''

Tôi thở dài một hơi nhẹ, nói tiếp: ``Ở cái thế giới nơi mà hôm nay bình yên nhưng ngày mai văn phòng có thể bị thổi bay, hay những sự cố bất thường có thể ập đến bất cứ lúc nào, thì tích cóp nhiều hơn chưa bao giờ là thừa cả. Em cần một nguồn tài chính đủ dày để bảo vệ gia đình mình trước mọi biến cố bất ngờ.''

``\dots Quả đúng là những gì mà một tên cuồng công việc sẽ thốt ra,'' giọng nói từ chiếc đồng hồ vang lên, không giấu nổi sự bỡn cợt, nhưng cũng có chút nghiêm túc.

Ryuuhou ngồi bên cạnh bĩu môi, quay sang nhìn tôi với vẻ mặt vừa thương vừa giận: ``Anh lại bắt đầu cái suy nghĩ cẩn trọng quá mức rồi đấy, Onii! Bọn em biết anh lo cho gia đình, nhưng lúc nào anh cũng chuẩn bị cho kịch bản tận thế như vậy thì làm sao mà sống thư thái nổi? Gia đình mình hiện tại đã ổn định hơn trước nhiều rồi mà!''

``Anh không cho rằng mình cẩn trọng quá mức đâu, Ryuuhou,'' tôi bình tĩnh cắt lời cô em gái. ``Sự chuẩn bị chưa bao giờ là thừa thải. Đến khi có chuyện thật sự xảy ra, tiền tích cóp chính là thứ duy nhất giúp chúng ta tự chủ được mọi thứ.''

A-san nhìn tôi một lúc rồi khẽ thở dài, nhẹ nhàng đẩy lại gọng kính trên sống mũi: ``Chuyện tích cóp phòng xa thì chị hiểu, nhưng em cũng thấy đấy, quy mô của COVER đang phát triển với tốc độ chóng mặt. Số lượng các thành viên tăng lên từng đợt, các dự án lớn nhỏ dồn dập đổ về. Chính em cũng biết một mình em không thể mãi gánh vác toàn bộ khối lượng công việc vận hành hậu cần này được, đúng không?''

Tôi lặng người. Đúng là vậy. Mấy tháng gần đây, dù đã cố gắng gồng gánh hết sức, tôi cũng nhận ra giới hạn sức lực của một con người. Công ty vừa mới mở thêm hai nhánh quốc tế mới là nhánh Indonesia và tiếng Anh, còn \hls\ - nhánh chỉ toàn nam giới mà Kaigou sắp tới sẽ đảm nhận vị trí tổng quản - nếu không có cô ấy gánh vác phía đó, thì mớ công việc vận hành kia cũng sẽ dồn hết lên vai tôi. 

Nếu quy mô cứ tiếp tục mở rộng với tốc độ này, dù có phân thân ra mười mảnh tôi cũng chẳng thể nào kham nổi tất cả mọi việc.

``Chị nói cũng có lý\dots\ Nghĩ lại thì, mấy tháng gần đây công việc ngập đầu thật, em làm mãi cũng không xuể.''

Chị đồng nghiệp mỉm cười, rút ra một bản đề xuất mới từ dưới xấp hồ sơ rồi đẩy về phía tôi: ``Chính vì thế, ban quản lý đã bàn bạc và đưa ra phương án mới. Bắt đầu từ mùa đông năm nay, công ty sẽ tuyển thêm nhân sự vận hành mới. Thay vì phải trực tiếp chạy đôn chạy đáo xử lý từng việc lớn nhỏ như trước, nhiệm vụ chính của em sẽ là hướng dẫn, đào tạo và giám sát các nhân viên đó làm việc. Em sẽ đóng vai trò điều phối toàn bộ hệ thống.''

Điều đó đồng nghĩa với việc khối lượng công việc chân tay và những đêm thức trắng liên miên của tôi sẽ giảm đi một cách đáng kể.

``Thật sao chị A?!'' Ryuuhou lập tức tròn mắt, hai tay đập nhẹ xuống mặt bàn, gương mặt sáng rực lên. Cô nàng quay sang lay mạnh vai tôi giục giã: ``Đồng ý đi Onii! Cơ hội tốt như thế này còn chờ gì nữa!''

Tôi hơi chần chừ. Mặc dù lời đề nghị nghe cực kỳ hấp dẫn, nhưng trong đầu một kẻ luôn ám ảnh về tài chính như tôi liền nảy ra một mối lo khác. Tôi ngập ngừng cất tiếng: ``Nhưng mà\dots\ chị A này. Chuyển sang làm người hướng dẫn, công việc nhàn hơn nhiều\dots\ thì tiền lương của em có bị ảnh hưởng hay cắt giảm không?''

Chị tiền bối khựng lại một giây rồi bật cười thành tiếng, lắc đầu thán phục trước sự thực dụng đến mức đáng quan ngại của tôi: ``Em đúng là\dots\ lúc nào cũng chỉ nghĩ đến tiền! Yên tâm đi Kuon-san, công ty trả lương cho em không phải vì sức lao động chân tay hay số giờ em thức đêm, mà là vì kiến thức, kinh nghiệm và khả năng xử lý khủng hoảng duy nhất chỉ có ở em. Việc em truyền dạy quy trình đó cho nhân viên mới chính là nhân rộng giá trị cho công ty. Mức lương cố định của em vẫn giữ nguyên không thiếu một xu, thậm chí còn có thêm phụ cấp trách nhiệm đào tạo đấy!''

Ngay lúc đó, chiếc đồng hồ trên cổ tay tôi phát ra tiếng bíp nhẹ. Giọng nói mỉa mai, điềm nhiên của Game lập tức vang lên ngắt lời: ``Cậu có vẻ lo xa quá, Kuon-user ạ, còn hơn cả mấy sinh viên xa nhà nữa. Nhưng mà đấy, Yuujin-san nói cũng có lý, chuyển từ `tự bóc phốt sức lao động' sang `ủy quyền cho tiến trình phụ' là quyết định tối ưu hóa tài nguyên hoàn hảo. Tôi nghĩ cậu nên chấp thuận ngay lập tức để giảm nguy cơ đột quỵ do thiếu ngủ.''

Tiếng cười vỡ ra từ cả hai phía khi trợ lý ảo của tôi lại một lần nữa bóc phốt tôi không chút nương tay. 

Tôi giả vờ trầm ngâm thêm vài giây, trong đầu đã dàn dựng sẵn cả một kịch bản. Nếu là cái bản ngã chỉ biết làm việc và tích tiền của mấy năm trước, có lẽ tôi đã lao vào nhận lời ngay lập tức khi nghe lương không giảm, công việc lại nhàn đi. Một lời đề nghị hoàn hảo không thể nào tuyệt vời hơn.

``\dots Em xin lỗi A-san, nhưng có lẽ em sẽ từ chối.''

Lời tôi vừa dứt, cả căn phòng bỗng chốc im lặng. A-san đẩy gọng kính lên, tròn mắt nhìn tôi không thể tin vào tai mình, còn Ryuuhou thì há hốc miệng, kéo mạnh cánh tay tôi như muốn đánh thức tôi khỏi cơn mê: ``Onii!? Anh đang nói gì thế!? Bộ cả đời này anh định bán lưng cho bọn tư bản tới chết sao!?''

Ngay chiếc đồng hồ trên cổ tôi cũng lập tức bật sáng, giọng Game đầy hoang mang cất lên: ``Cậu bị điên rồi sao, Kuon-user!? Một cơ hội ngon ăn thế này mà cậu lại từ chối là thế nào?''

Tôi vội giơ tay lên, ra dấu cho cả hai người giữ bình tĩnh lại. Tôi khẽ cười một lúc, rồi nói: ``\dots Đó là những điều mà em sẽ nói nếu em còn là một kẻ chỉ biết đắm chìm trong công việc và tiền bạc như trước. Nhưng trong suốt chặng đường chu du qua các vũ trụ khác nhau, em đã tự hỏi mình rất nhiều lần: rốt cuộc, tiền bạc có phải là tất cả của cuộc đời này?''

Thật vậy, hồi còn mắc kẹt ở thế giới ERROR, trong những ngày hiếm hoi mà chúng tôi không phải chạy đua với vòng lặp chết chóc hay vật lộn để sinh tồn, tôi đã có cơ hội ngồi lại và thực sự nhìn nhận lại cuộc đời mình. Lúc đó, giữa thị trấn mịt mù sương và những lo lắng chỉ xoay quanh việc làm sao để tồn tại qua ngày hôm sau, tôi chợt nhận ra rằng tiền bạc hay vật chất chẳng là gì cả. Điều thực sự quan trọng không phải là chúng ta tích lũy được bao nhiêu, mà là chúng ta đã sống như thế nào, đã trân trọng được bao nhiêu khoảnh khắc bên người thân yêu.

Ryuuhou chớp chớp mắt nhìn tôi, giọng nhỏ lại: ``Vậy thì anh nhận ra được gì rồi, Onii?''

``Tức là, điều quan trọng nhất khi em làm ở đây không phải là có thể đưa về cho gia đình bao nhiêu tiền, mà là việc em có tận hưởng ở đây hay không,'' tôi tiếp tục, trong đầu hiện lên những kỷ niệm vui buồn cùng các cô gái ở Tam Giác Xanh, những chuyến đi chơi, những bữa ăn tối ấm áp, cả những lúc chúng tôi cùng nhau vượt qua khó khăn. 

``Trước đây em cứ nghĩ rằng phải gánh vác mọi thứ một mình, phải làm việc đến kiệt sức mới là có trách nhiệm. Nhưng bây giờ em đã hiểu, để có thể thực sự tốt với gia đình, trước hết em phải tốt với chính mình. Để có thể chăm sóc mọi người, em phải biết chăm sóc sức khỏe của bản thân trước. Chia sẻ công việc không phải là trốn tránh trách nhiệm, mà là để chúng ta cùng nhau phát triển, để em có thêm thời gian bên gia đình, cùng các em tạo thêm nhiều kỷ niệm hơn nữa.''

``Vậy, có nghĩa là\dots'' A-san vẫn còn ngỡ ngàng, chỉ có thể lắp bắp hỏi lại.

``Em đồng ý,'' tôi thở phào nhẹ nhõm, gật đầu chắc nịch với chị đồng nghiệp của mình.

``\textbf{Tuyệt vời!}'' Ryuuhou reo lên đầy vui sướng. Cô em gái nhảy chồm tới ôm chặt lấy cánh tay tôi, cười rạng rỡ: ``Thế là từ giờ em sẽ không phải thấy cảnh Onii thức xuyên đêm xanh xao như bóng ma nữa rồi! Anh cũng sẽ có thêm thời gian để tập trung hoàn thành nốt mấy cái đồ án đại học gắt gao trên trường nữa chứ!''

Chiếc đồng hồ lại bật lên tiếng bíp, giọng Game lần này tràn đầy sự hài lòng: ``Cuối cùng cũng đưa ra quyết định đúng đắn, Kuon-user. Hy vọng từ nay tỉ lệ ngủ ít hơn 5 tiếng mỗi ngày của cậu sẽ giảm xuống, để tôi không phải thường xuyên thông báo nguy cơ đột quỵ nữa.''

Em tôi cũng nối ngay: ``Đúng đó Game-san! Từ nay chúng ta có thể cùng nhau đi chơi nhiều hơn nữa, không phải cứ ở nhà chờ Onii làm việc nữa!''

Tôi đứng đó, nhìn nụ cười rạng rỡ của em gái, nghe tiếng cười thoải mái của A-san và giọng nói vẫn hay trêu chọc của Game, trong lòng thấy ấm áp lạ thường. Đúng là cuộc sống có thật nhiều điều quý giá hơn tiền bạc, và tôi thật may mắn khi đã kịp thời nhận ra điều đó, để có thể bắt đầu một chương mới thật sự bình yên và hạnh phúc bên gia đình nhỏ của mình.

\ct

Sau lúc cả ba cùng cười đùa ăn mừng chút xíu ấy, A-san khẽ nở một nụ cười nhẹ, bỗng dưng như sực nhớ ra chuyện gì, chị liền xoay người nhìn sang Ryuuhou với ánh mắt thâm thúy, đầy ẩn ý.

``Nhân tiện nói mới nhớ, Ryuuhou-san này\dots'' A-san chống cằm, nghiêng đầu hỏi nhỏ. ``Mấy tháng nay chị để ý thấy em đến văn phòng công ty ngày càng thường xuyên quá nhỉ. Đặc biệt là trong ba tháng qua, gần như ngày nào chị cũng thấy em lượn lờ ở đây.''

Ryuuhou hơi ngẩn người một chút, đưa tay gãi gãi mảng tóc ngắn của mình: ``Hả? Thì\dots\ em lên chơi với Onii thôi mà chị, tiện thể phụ anh ấy mấy việc vặt vãnh trong văn phòng thôi.''

``Thế nên chị mới nghĩ ra ý tưởng này đấy,'' chị đồng nghiệp mỉm cười gật đầu. ``Công ty sắp tới có nhiều đợt ra mắt thành viên mới và các dự án biểu diễn sân khấu, nên đang rất cần người phụ trách chuyên môn về mảng luyện thanh, biên đạo nhảy và sắp xếp vị trí đội hình cho các thành viên. Mấy việc này vốn dĩ em đã làm rất tốt với vai trò trưởng nhóm nhạc underground của em rồi đúng không?''

``Dạ\dots\ vâng?'' con bé nghiêng đầu đáp lại. ``Có chuyện gì hả chị?''

``Em có muốn chính thức đến thử việc ở COVER với vị trí này không?''

Vừa nghe nhắc đến lĩnh vực mà mình đam mê, đôi mắt của Ryuuhou lập tức long lanh sáng rực. Cô em gái vỗ mạnh vào ngực cái *BỤP*, trả lời đầy tự tin mà không hề do dự: ``Chuyện đó chị hoàn toàn yên tâm giao phó cho em! Mấy vụ luyện giọng, chia bè hay xếp vị trí trên sân khấu sao cho thật bùng nổ thì em làm đã quá quen tay rồi. Em đồng ý ngay! Ngay cả khi bắt đầu thử việc từ tuần sau em cũng hoàn toàn ổn ạ!''

Nhìn thấy cô em gái hào hứng nhận lời nhanh như chớp, tôi giật mình thon thót, vội vàng lên tiếng ngăn lại: ``Chờ đã nào Ryuuhou! Chuyện này đâu phải trò đùa đâu. Em đã suy nghĩ thật kỹ chưa? Còn lịch tập, lịch thu âm của ban nhạc underground bên em thì sao? Lỡ nhận thêm việc ở COVER mà khối lượng công việc dồn dập, ảnh hưởng đến hoạt động của nhóm thì em giải quyết ra sao?''

Tôi thực sự đứng ngồi không yên. Ban nhạc là đam mê lớn nhất đời của Ryuuhou, tôi không muốn con bé vì thấy anh trai vất vả mà lại đặt ước mơ của mình sang một bên, hay ép bản thân phải làm việc quá sức.

Trái ngược với sự lo lắng của tôi, con bé quay sang nhìn tôi, trong mắt chứa đầy sự kiên định và chân thành. Ryuuhou mỉm cười nhẹ nhàng, như muốn xoa dịu nỗi lo của tôi: ``Onii lo xa quá rồi! Lịch tập của nhóm em vốn rất linh hoạt, toàn bộ xếp vào buổi chiều tối hoặc cuối tuần thôi, em hoàn toàn có thể cân bằng được cả hai công việc. Thêm vào đó\dots\ miễn là công việc này có thể giúp được Onii, san sẻ bớt gánh nặng cho anh thì chút vất vả nhỏ đó em có ngại gì đâu!''

Lời nói ấm áp của cô em gái khiến tim tôi như hẫng đi một nhịp, bao nhiêu lời căn nhắc định nói ra đành phải nuốt ngược vào trong bụng.

``Nghe kiểu gì cô bé cũng làm tất cả những việc này vì ông anh cả của mình nhỉ?'' giọng Game lại vang lên với giọng điệu châm biếm từ chiếc đồng hồ trên cổ tay tôi. ``Tuy nhiên, việc bổ sung thêm một Ryuuhou-user vừa có chuyên môn cao lại còn có chi phí nhân sự `người nhà' thì chính là quyết định tốt nhất cho tình hình tài chính của cậu đó, Kuon-user.''

Nghe ông ta nói vậy, em gái tôi liền vỗ tay liên tục, gương mặt rạng rỡ như vừa nhận được món quà lớn nhất đời: ``Thế là đã quyết định rồi nhé! Từ tuần sau em sẽ chính thức đến COVER làm việc, cùng Onii và mọi người cùng nhau cố gắng!''

Tôi đứng đó nhìn nụ cười tươi như hoa của em gái, trong lòng tràn đầy một sự ấm áp khó tả. Không ngờ rằng chỉ trong một buổi sáng ngắn ngủi, cả tôi và em đều có thêm một bước ngoặt mới trong công việc. Có thêm em ở bên cạnh san sẻ gánh nặng, tôi tin rằng mọi chuyện sau này sẽ đều diễn ra thuận lợi. Tôi cũng có thể dành nhiều thời gian hơn cho gia đình, tận hưởng những khoảnh khắc bình yên bên những người mình thương yêu, thay vì cứ mãi chìm đắm trong đống công việc không hồi kết.

\ct

Sau khi đã xong xuôi chuyện của Ryuuhou, A-san nhẹ nhàng bắt chéo hai tay, gõ nhẹ các ngón tay lên mặt bàn với vẻ mặt trầm ngâm, có chút băn khoăn. Cô ngước nhìn tôi, thở dài từ từ: ``Nhắc đến nhân sự trong gia đình em\dots\ chị vẫn còn một việc khá là lo lắng. Đó là về vị trí tổng quản \hls\ của Kaigou-san.''

``Kaigou ạ?'' tôi nhíu mày lại, bất ngờ.

``Ừm,'' chị ấy gật đầu, ánh mắt đầy ái ngại. ``Em cũng rõ cái tính ghét đàn ông như hắt nước đổ đi của con bé rồi mà. Lần đầu tiên chị và sếp Tanigo-san ngỏ lời mời cô ấy nhận vị trí quản lý cho nhánh nam \hls, con bé đã một mực từ chối ngay lập tức, thậm chí còn nhìn bọn chị như thể vừa đưa ra một đề nghị điên rồ nhất trần gian vậy.''

Ngay lập tức, đầu tôi liền nhức nhối hồi tưởng lại chuỗi sự kiện trước đây. Từ cái hồi còn ở The Void, Kaigou suýt nữa thì đấm chết tôi  - lúc đó cũng có thể thông cảm, vì cô nàng chưa quen biết tôi và vốn đã giữ thái độ đề cao cảnh giác với toàn bộ cánh đàn ông, rồi đến lúc ở Nijisanji, cô ấy coi các Livers nam như chuột bọ, luôn tránh né xa. Đến cả khi ở Aogami, cô nàng cũng từng sẵn sàng xông vào đánh túi bụi đám con trai dám bắt nạt Suzuno của mình. Suốt bao nhiêu lần như thế, tôi chỉ có thể đứng chết trân nhìn cô em song sinh ``hành sự'' mà chẳng thể làm gì được. Một cô gái với cái tính cách ``trời đánh'' cùng tư duy giao tiếp lệch lạc thế này, nếu không được uốn nắn đàng hoàng thì làm sao có thể đảm đương nổi vị trí tổng quản ở \hls\ nơi mà chỉ có nam giới được đây?

Nghe đến đó, tôi chỉ biết dựa lưng vào ghế và thở một hơi dài não nề.

``Chị nói đúng lắm, A-san. Kaigou nhà em vẫn còn phải dạy bảo nhiều. Cái thù hằn với nam giới từ những thế giới cũ khiến em ấy lúc nào cũng xù lông nhím, sẵn sàng phản ứng bất cứ khi nào phải tiếp xúc với con trai.''

Nhớ lại những chuyện xưa cũ, tôi lắc đầu ngao ngán rồi tiếp tục kể: ``Chị có biết cái ngày Kaigou chính thức hoàn tất thủ tục nhập học và bước chân vào lớp đại học của em không? Dù bị chậm mất tận năm học kỳ - tương đương hai năm rưỡi - vì mấy rắc rối về thủ tục trú ngụ và thời gian phục hồi thể trạng, nhưng cô nàng vừa xuất hiện đã làm cả lớp một phen toát mồ hôi hột.''

\begin{center}
    \uline{Hồi tưởng\\Thứ hai vừa rồi\\Phòng 312-H1 - Đại học Xây dựng Kawachi.}
\end{center}

Tôi nhớ như in buổi sáng hôm đó. Lớp CS của tôi có hơn bốn mươi tên con trai và vỏn vẹn đúng bốn bạn nữ. Khi Kaigou bước vào, áo khoác đen khoác hờ trên vai, đôi mắt sắc lẹm lướt qua từng gò má nam sinh như thể đang nhìn một mớ rác thải công nghiệp cần phân loại gấp. 

``\vie{Chào. Hikawa Kaigou đây.}'' Cô ấy chỉ đáp cụt lủn như vậy, mắt thậm chí còn chẳng buồn liếc về phía tôi, trông không khác gì một tên trùm trường xã hội đen.

Áp lực vô hình từ sát khí của cô nàng khiến cả dãy bàn đầu đồng loạt gục mặt đổ mồ hôi hột. Thậm chí, tới giảng viên bộ môn của lớp còn phải cũng phải lúng túng mất vài giây mới lấy lại được bình tĩnh. Bốn bạn nữ trong lớp - Yoisai trong ``nhóm đồ án'' của chúng tôi là một trong số đó - những người duy nhất nhận được một cái gật đầu nhẹ nhàng cùng ánh mắt dịu bớt đôi chút từ cô. 

``\vie{\hl{N-Này, Kuon ơi\dots{} ai vậy? Sao cô ấy toát ra thứ\dots{} sát khí đáng sợ thế?}}'' Tamachi ngồi bên cạnh, ghé sát vào tai tôi hỏi nhỏ.

``\vie{\hl{\dots Có lẽ các ông không tin, nhưng đó là em gái tôi đấy,}}'' tôi tặc lưỡi đáp lại, nhìn nét mặt đang căm căm của cô nàng song sinh kia.

``\vie{\hl{Ông còn có cả em gái song sinh nữ á?}}'' Yoisai ở bàn dưới ngạc nhiên đáp lại. ``\vie{\hl{Tôi tưởng nhà ông chỉ có Ryuuhou là em gái ông thôi chứ?}}''

``\vie{\hl{\dots Chuyện dài lắm,}}'' tôi đáp gọn, ``\vie{\hl{Không tiện kể đâu.}}''

``\vie{\hl{S-Sao em gái c-của ông lại đ-đáng sợ thế\dots}}'' Yamazu ở phía bên kia bàn cũng toát mồ hôi hột, tay chân hơi bủn rủn.

Thấy không khí quá căng thẳng, tôi đành nhổm người dậy tính cất lời chữa cháy: ``Kaigou à, em--''

``\textbf{Im đi!}''

Cô quay sang trừng mắt quát thẳng vào mặt tôi với âm lượng đủ làm rung chuyển cái bảng đen. Tôi đứng hình mất ba giây, còn đám con trai xung quanh thì run rẩy co rúm lại vì tưởng tôi sắp bị cô nàng hạ nốc-ao tại chỗ, không dám ho he lấy một tiếng.

Rốt cuộc, cả buổi học hôm đó diễn ra trong sự ngột ngạt đến nghẹt thở, tới giảng viên thi thoảng cũng phải liếc qua cô học sinh mới, nơm nớp không biết liệu thầy ấy có làm phật lòng không, và vì thế, chẳng một ai trong lớp có thể tập trung nổi vào bài giảng. 

Vừa khi tiếng chuông báo hết tiết học vang lên, đám nam sinh lập tức cuống cuồng gấp gọn sách vở rồi xô nhau chạy ra khỏi phòng, y hệt như bọn người vừa may mắn thoát thân khỏi hang ổ của yêu quái. 

``\vie{Trời ơi\dots\ mới tưởng sắp bị cô nàng kia ăn tươi nuốt sống mất!}'' Tamachi thở phào đầy sợ hãi, vừa chạy vừa vụng về xách theo chiếc ba lô của mình. ``\vie{Em gái ông đáng sợ quá!}''

``\vie{Đ-Đúng đấy! Ánh mắt cô ấy s-s-ắc lẹm thật, tôi mà đứng thêm năm giây nữa ch-chắc tim ngừng đập mất!}'' Yamazu nối lời, hai chân vẫn chưa hết run rẩy khi vừa thoát khỏi không khí ngột ngạt trong lớp.

Yoisai cũng đi cạnh họ, quay đầu nhìn lại phía phòng học rồi lắc đầu cười khổ: ``\vie{Từ trước đến giờ tôi chưa gặp ai mà khiến cả một lớp bốn mươi thằng con trai sợ hãi đến thế đâu. Anh Kuon mà còn không đứng ra dàn xếp, không biết có chuyện gì xảy ra nữa.}''

Tôi vừa chậm rãi xếp lại quyển giáo trình của mình, vừa đáp lại với tiếng thở dài não nề: ``\vie{Mọi người cứ về trước đi. Có gì tôi chịu trận thay cho.}''

``\vie{B-Bảo trọng!}'' cậu bạn Đạo diễn đáp lại, rồi phóng thẳng ra ngoài hành lang, cố gắng chạy càng xa càng tốt. Cô bạn Lập trình và cậu bạn Đeo kính cũng nối gót theo.

Khi bóng dáng của người cuối cùng đã biến mất hẳn ở cuối hành lang, Kaigou mới khoanh tay trước ngực, quay lưng về phía tôi rồi gầm gừ với một giọng trầm thấp, ``\vie{Hikawa Kuon. Ở lại đây. Tôi có chuyện cần nói.}''

Tôi thở dài xếp lại xấp tài liệu. Thế nhưng, ngay khi tiếng bước chân ngoài hành lang hoàn toàn dứt hẳn, cái dáng vẻ ngầu lòi, lạnh băng ban nãy của cô nàng đột ngột sụp đổ hoàn toàn.

Cô em gái song sinh quay ngoắt lại, hai tay cuống quýt xua xua, gương mặt đỏ bừng vì hoảng hốt và xấu hổ. Cô lao nhanh về phía bàn tôi, gập người chín mươi độ liên tục như một chiếc máy dập: ``\vie{Em xin lỗi em xin lỗi em xin lỗi! Em thật sự xin lỗi anh!}''

``\vie{\dots Cứ bình tĩnh đã nào,}'' tôi xoa xoa thái dương. ``\vỉe{Anh hiểu là em thực sự ghét đàn ông và anh cũng lường trước chuyện này rồi, cơ mà\dots}''

``\vie{Lúc nãy em không cố ý cao giọng với anh đâu!}'' Cô nàng lí nhí, hai ngón tay trỏ vô thức chọc chọc vào nhau, mắt rấn rấn nước vì hối hận. ``\vie{Tại\dots\ tại nhìn chung quy toàn đực rựa làm em bị ngợp quá, phản xạ tự nhiên nó tự bộc phát\dots\ Anh không giận em chứ?}''

``\vie{Không,}'' tôi chỉ đáp gọn như thế. Nhìn cô em gái vừa phút trước còn như ngọn núi lửa chực chờ phun trào, giờ đây lại xoắn xuýt hối lỗi đến mức tội nghiệp chỉ vì lỡ gằn giọng với anh mình, tôi chỉ biết lắc đầu ngao ngán.

\begin{center}
    \uline{Kết thúc hồi tưởng.}
\end{center}

Quay về với hiện tại, tôi chắp tay sau đầu nhìn A-san: ``Đấy, đến việc cư xử bình thường trong một lớp học toàn con trai mà cô ấy còn cuống cuồng mất phương hướng như thế, huống hồ gì làm quản lý \hls\ mà không cảm thấy áp lực cho được.''

``Cũng không thể trách Kaigou-nee được đâu,'' Ryuuhou chen lời, hai tay đan vào nhau đặt trên bàn. ``Nhưng mà tới mức làm nguyên cả lớp của Onii sợ sệt như vậy\dots\ rồi xin lỗi rối rít anh trai mình, em mừng tượng mà thấy buồn cười.''

Chiếc đồng hồ trên cổ tay tôi cũng bật sáng, giọng Game bâng quơ xen vào: ``Đúng vậy, Ryuuhou-user nói không sai. Từ số liệu hành vi mà tôi thu thập được trong khoảng thời gian vừa qua, Kaigou-user có đến ba mươi lần tiếp xúc không cần thiết với đối tượng nam giới lạ đều dẫn đến hành vi tránh né hoặc thể hiện thái độ thù địch rõ rệt. Thử đặt cô ấy vào một môi trường toàn nam giới như \hls\ mà không có ai quen ở bên, nguy cơ xảy ra xung đột là rất cao.''

Chị đồng nghiệp nghe hết câu chuyện của tôi mà không thể kiềm được tiếng cười khúc khích. Cô gật gù liên tục, rồi ánh mắt dần trở nên nghiêm túc, lấp đầy niềm tin vững chắc vào tôi.

``Ra là vậy\dots\ Thế thì chị càng có lý do chính đáng để giao phó trọng trách này cho em rồi, Kuon-san. Ngay từ khi bắt đầu công việc tại \hls, em hãy trực tiếp đảm nhận việc hướng dẫn Kaigou-san. Không chỉ đơn thuần là truyền đạt lối làm việc chuyên nghiệp, mà em còn hãy giúp cô ấy cởi mở hơn, định hình lại cách ứng xử cũng như thay đổi cái nhìn thận trọng quá mức đối với các bạn nam.''

Tôi ngạc nhiên nhướng mày, liền cất tiếng hỏi: ``Hả? Em tưởng chị A đã đồng ý sẽ tự mình hướng dẫn em ấy mà? Chuyện lớn thế này sao lại đổ hết lên vai em?''

Game bật cười khanh khách từ chiếc đồng hồ, giọng đầy giễu cợt: ``Cậu còn hỏi gì nữa, Kuon-user? Xét trên phương diện thống kê, cậu là người duy nhất trong tất cả các đối tượng nam giới mà Kaigou-user không hề có bất kỳ hành vi phản kháng nào, thậm chí còn thể hiện sự tin tưởng tuyệt đối. Giao cậu làm người hướng dẫn là lựa chọn hợp lý nhất mà ban quản lý có thể đưa ra rồi đấy.''

A-san liền khẽ cười, vắt chéo chân và giải thích thêm: ``Game-san nói đúng đấy. Chị vốn định như vậy, nhưng sau khi nghe em kể lại chuyện ở lớp đại học mới biết tình hình phức tạp hơn chị nghĩ. Chị không hiểu rõ gốc gác cũng như tính cách của Kaigou-san bằng em, nếu chị tự ý hướng dẫn có khi còn khiến con bé căng thẳng thêm, thậm chí còn phản tác dụng nữa. Hơn nữa, ban quản lý cũng từng đề cập rằng có em, người thân thiết nhất với Kaigou-san, đứng ra kèm cặp thì sẽ hiệu quả hơn rất nhiều.''

``Nói tóm lại, chỉ có Onii mới có thể làm được điều đó thôi,'' Ryuuhou ngồi bên cạnh tóm gọn lại, vỗ tay khe khẽ như thể vừa nhấn mạnh một sự thật hiển nhiên. ``Thử để người khác dạy dỗ Kaigou-nee xem, chưa kịp nói hết câu đã bị chị ấy dập cho ngã ngửa rồi.''

Tôi thở dài bất lực, đưa tay gãi gãi má. Quả nhiên, trọng trách nặng nề này lại một lần nữa đổ hết lên vai tôi. Nhìn ánh mắt tin tưởng của cả A-san lẫn Ryuuhou, còn giọng Game vẫn liên tục nhắc nhở rằng đây là ``quyết định tối ưu hóa rủi ro duy nhất', tôi cũng không thể nào từ chối được nữa, chỉ đành gật đầu nhận lời.

\ct

``À, đúng rồi, suýt nữa thì em quên\dots'' Tôi đưa tay gõ nhẹ lên xấp báo cáo thống kê thiệt hại trên bàn: ``Nhắc đến vụ các thành viên phá hoại văn phòng\dots\ Em cũng vừa hoàn thiện một giải pháp triệt để cho vấn đề này. Bắt đầu từ tháng sau, chị sẽ không còn phải đau đầu nhìn mấy con số đền bù thiết bị hay sửa chữa cơ sở vật chất nữa đâu.''

A-san chớp mắt, tỏ vẻ tò mò: ``Giải pháp gì cơ?''

``Một thiết bị tự động khôi phục công trình,'' tôi mỉm cười giải thích. Nó được tích hợp mạch quét vi mô và cơ chế tái tạo vật chất. Cứ mỗi khi các thành viên lỡ tay bộc phát năng lượng làm đập phá trần nhà, đập gãy bàn ghế hay làm vỡ cửa kính, hệ thống sẽ lập tức thu gom các mảnh vỡ, tự động lắp ráp và hoàn nguyên mọi thứ về trạng thái ban đầu trong vài giây. Vừa tiết kiệm tối đa chi phí sửa chữa cho công ty, lại vừa cho các cô gái thỏa thích\dots\ phá tung cả văn phòng mà không lo để lại hậu quả.''

Nghe tới đây, chị đồng nghiệp khẽ chau mày, ánh mắt hiện rõ sự e dè và hoài nghi: ``Tự động hoàn nguyên sao? Nghe thì tuyệt đấy, nhưng công nghệ này có thực sự ổn định không? Chị thấy hơi lo lo đấy\dots\ Nhỡ đâu nó lại chập điện hay bùng nổ thì sao?''

``Để em demo thực tế cho chị xem nhé,'' tôi quay sang nhìn cô em gái bên cạnh, hất cằm về phía ô cửa sổ kính cường lực lớn hướng ra ban công. ``Ryuuhou, làm mẫu một phát cho chị ấy xem nào.''

``Nhận lệnh, Onii!'' con bé khoái trí reo lên. Cô nàng chẳng hề ngần ngại, nhún người một cái rồi tung một cú đá uy lực thẳng vào ô cửa sổ kính ngàn đô của văn phòng.

*XOẢNG!*

Tiếng kính vỡ tan tành vang lên chói tai. Mảnh kính bắn tung tóe khắp nơi, gió lạnh từ bên ngoài tràn ào ạt vào phòng.

Chị đồng nghiệp đeo kính giật mình bật đứng dậy, hai tay ôm chặt lấy đầu, gương mặt nhợt nhạt không còn chút màu sắc. Cô vừa hoảng loạn vừa kinh ngạc kêu lên: ``Trời ơi! Ryuuhou-san! Em làm gì thế này!? Cái cửa kính đó--''

Nhưng trước khi chị ấy kịp nói hết câu trách móc, một luồng ánh sáng ma trận xanh lam mỏng manh từ thiết bị gắn trên tường lập tức quét qua toàn bộ khu vực khung cửa. Những mảnh kính vụn đang vương vãi khắp sàn nhà bỗng nhiên bay lên trời, tự động xoay chuyển và lắp ráp lại với nhau với tốc độ chóng mặt, chẳng khác nào đang tua ngược lại một đoạn phim. Chỉ trong nháy mắt, ô cửa sổ lớn đã lành lặn hoàn toàn, như thể chưa từng có chuyện gì xảy ra: một ô cửa sáng bóng, nhẵn nhụi, chẳng có lấy một vết xước nhỏ. Gió lạnh từ bên ngoài lập tức bị chặn lại, căn phòng nhanh chóng quay trở lại trạng thái yên tĩnh ban đầu.

A-san đứng chết trân giữa phòng, hai mắt mở to hết cỡ, gương mặt từ tái mét vì sốc chuyển dần sang bàng hoàng, kinh ngạc, rồi cuối cùng vỡ òa trong sự ngỡ ngàng không thể tả. Đôi mắt sau tròng kính sáng rực lên như vừa chứng kiến một phép màu hiện hữu trước mắt. Cô bước chậm rãi lại gần, bàn tay run rẩy vuốt nhẹ lên bề mặt kính phẳng lì, hoàn toàn nguyên vẹn.

``Không\dots\ không thể nào!'' chị ấy thốt lên, giọng nói thay đổi hoàn toàn 180 độ. ``Tất cả mảnh vỡ\dots\ tự động dính lại và phục hồi như mới hoàn toàn sao?!''

``Đã bảo là thuật toán tái cấu trúc của tôi không bao giờ có sai số rồi mà, thưa madame,'' giọng Game vang lên từ chiếc đồng hồ trên cổ tay tôi, chứa đựng chút tự hào ẩn giấu. ``Từ giờ trở đi, dù mọi người có biến cả văn phòng thành chiến trường thì chi phí bảo trì hằng tháng vẫn sẽ là\dots\ con số không tròn trịa.''

Tôi khoanh tay dựa lưng vào ghế, nở một nụ cười thỏa mãn nhìn vẻ mặt kinh ngạc của chị đồng nghiệp. Ryuuhou cũng vỗ tay cười khoái chí, chỉ vào tôi nói với A-san: ``Chị thấy chưa! Onii em tính toán hết cả chuyện đó từ lâu rồi ấy mà. Cái thiết bị tự động sửa chữa này được anh ấy chuẩn bị từ mấy tháng trước, chỉ đợi hôm nay mới dám demo cho chị xem thử hiệu năng thôi!''

\ct

Sau khi đã chốt xong các vấn đề chính liên quan đến nhân sự, buổi họp tiếp tục lan man sang những chủ đề bên lề, từ hiệu suất của các sự kiện live gần đây đến chuyện trêu chọc mấy thành viên mới vô tình làm hỏng trang thiết bị trong studio. A-san vừa lật lướt vài trang báo cáo nhỏ, vừa chia sẻ thêm vài tin tức vui vẻ trong công ty, còn Ryuuhou thì ngồi bên cạnh hay nhổm người lên hỏi thêm chuyện về các dự án biểu diễn sắp tới, đôi mắt càng lúc càng sáng rực khi nghe đến cơ hội được phụ trách mảng âm nhạc. Game cũng thỉnh thoảng xen vào vài câu bình luận mỉa mai từ chiếc đồng hồ, làm cả phòng cười nghiêng ngả khi cứ liên tục ``bóc phốt'' những sai lầm nhỏ nhặt của tôi trong các báo cáo tháng trước.

Cuộc họp kết thúc lúc mười hai giờ trưa. A-san gập gọn xấp hồ sơ báo cáo lại, nở một nụ cười nhẹ nhõm rồi thở phào một hơi dài: ``Vậy là buổi báo cáo tháng này kết thúc vô cùng tốt đẹp! Dù danh mục `thiệt hại' do các thành viên gây ra vẫn khiến người ta xót ruột như mọi khi, nhưng hiệu suất chung của em vẫn tuyệt đối hoàn hảo. Hơn cả thế, hôm nay chúng ta đã chốt được những bước chuyển giao nhân sự cực kỳ quan trọng: Kuon chuẩn bị chuyển sang vị trí điều phối và hướng dẫn, Ryuuhou-san sẽ chính thức thử việc mảng âm nhạc từ tuần sau, và kế hoạch `uốn nắn' Kaigou cho nhánh \hls\ cũng đã có hướng đi rõ ràng.''

Ryuuhou vừa vỗ tay khe khẽ, vừa hất cằm tự hào: ``Đúng vậy! Từ giờ em sẽ là đồng nghiệp chính thức của mọi người, không còn là người chỉ lượn lờ đến chơi nữa đâu!''

Game cũng lập tức bật lên tiếng cười khúc khích, giọng đầy trêu chọc: ``Thế thì Ryuuhou-user cũng phải chuẩn bị tinh thần mà làm việc chăm chỉ đấy nhé, đừng để đến lúc phải nhận lương mà chỉ biết đến ăn chơi nữa.''

``Dĩ nhiên!'' cô em gái tôi ưỡn ngực đáp lại.

Chị đồng nghiệp xoay ghế, hướng ánh mắt ra khung cửa kính lớn nhìn xuống thành phố Akiba tấp nập. Sự mệt mỏi trên gương mặt cô nhường chỗ cho sự hào hứng và kỳ vọng tràn trề: ``Giai đoạn tiếp theo của \hll\ sẽ là một chặng đường bùng nổ hơn bao giờ hết. Công ty đang tích cực chuẩn bị cho đợt ra mắt các thế hệ thành viên mới, lên kế hoạch cho những sự kiện đại nhạc hội Live hoành tráng quy mô lớn, đồng thời đẩy mạnh bước chân ra thị trường quốc tế. Mọi bánh răng đều đang quay với tốc độ chưa từng có. Sự đồng hành và san sẻ của ba anh em em lúc này chính là điểm tựa vững chắc nhất để COVER tự tin bước vào kỷ nguyên mới.''

Nhìn xấp tài liệu được sắp xếp gọn gàng trên bàn cùng gương mặt rạng rỡ, tràn đầy năng lượng của Ryuuhou bên cạnh, tôi cảm nhận được một luồng sinh khí mới đang len lỏi qua từng góc nhỏ của văn phòng.

Những giông bão từ sự kiện đa vũ trụ kia đã thực sự lùi xa. Cuộc sống bình thường mới không chỉ mang đến sự ổn định cho căn nhà nhỏ của gia đình Hikawa, mà còn mở ra một bước ngoặt lớn trong sự nghiệp của tôi. Tôi sẽ không còn phải cô độc quay cuồng trong những đêm thức trắng để gánh vác toàn bộ khối lượng công việc khổng lồ nữa. Giờ đây, phía sau tôi đã có sự hỗ trợ từ Ryuuhou, sự trưởng thành từng ngày của Kaigou, và một lộ trình rõ ràng để chuẩn bị cho tương lai.

``Như vậy là một khởi đầu mới khá tốt đẹp đấy nhỉ?'' giọng Game cất lên chậm rãi từ chiếc đồng hồ trên cổ tay, cắt ngang dòng suy nghĩ của tôi. ``Với những gì diễn ra như vậy, sao chúng ta không ăn mừng nhỉ?''

Ryuuhou bật cười, nắm lấy tay tôi kéo đứng dậy: ``Game-san nói đúng đấy! Đi ăn mừng thôi Onii! Tuần sau em chính thức đi làm rồi, anh phải khao em một bữa thật ngon đấy nhé! Chọn nhà hàng món Nhật mà em thích hồi trước đi, được không?''

Tôi mỉm cười, gật đầu chào A-chan rồi cùng cô em gái bước ra khỏi phòng làm việc. Ánh nắng ấm áp của ngày đầu tháng trải dài trên hành lang công ty. Chặng đường phía trước của \hll\ chắc chắn sẽ còn nhiều thử thách và biến số, nhưng với một hậu phương vững chắc và những người đồng hành đáng tin cậy, chúng tôi đã hoàn toàn sẵn sàng cho giai đoạn tiếp theo.

%==========================================TRUYỆN PHỤ=========================================
\omk

\lang{Etsunan}

Rời khỏi tòa nhà văn phòng, cái không khí dịu nhẹ và se lạnh của buổi trưa đầu tháng như gột rửa sạch sẽ mọi mệt mỏi sau buổi sáng làm việc kéo dài. Tôi chậm rãi bước đi trên vỉa hè, tâm trí vẫn còn tựa như đang nhâm nhi lại toàn bộ kết quả của buổi báo cáo sáng nay. 

Phải công nhận, mọi thứ diễn ra êm đẹp ngoài sức tưởng tượng. Không chỉ trút bỏ được gánh nặng công việc chân tay bấy lâu để chuyển sang vai trò người điều phối, tôi còn giải quyết triệt để mối lo chi phí sửa chữa bằng thiết bị khôi phục tự động. Hơn hết, việc định hướng lại vị trí cho Kaigou và đón nhận thêm Ryuuhou vào đại gia đình nhân sự công ty đã vẽ ra một viễn cảnh vô cùng yên tâm cho chặng đường sắp tới.

Đi bên cạnh tôi, cô em gái dường như chẳng giấu nổi sự phấn khích. Con bé vừa bước đi vừa nhảy chân sáo tung tăng, miệng ngân nga một giai điệu vui tươi không thành lời. Việc chính thức được nhận vào thử việc mảng luyện giọng và biên đạo, nghĩa là từ giờ sẽ được kề vai sát cánh làm việc cùng anh trai, và có thể là cả Kaigou ở nhánh \hls\ nữa, khiến năng lượng trong người cô nàng bùng nổ hơn bao giờ hết.

Chỉ có điều\dots\ con bé dường như quên mất hôm nay mình đang mặc cái gì.

Thấy mỗi cú bật nhảy của cô em gái lại làm tà váy denim ngắn nẩy lên bần bật theo nhịp bước, tới độ tôi có thể nhìn thấy một mảnh vải cotton màu vàng cam lấp ló bên trong, tôi không khỏi hắng giọng, cất tiếng nhắc nhở bằng một giọng điệu vô cùng thản nhiên: ``Ryuuhou này\dots\ Em nhảy chân sáo hăng hái quá đấy. Hàng họ của em lộ hết ra ngoài rồi kìa.''

Ryuuhou lập tức khựng lại giữa chừng. Hai gò má con bé đỏ bừng lên như quả cà chua chín, đôi phồng má giận dỗi quay ngoắt sang nhìn tôi. Ngay lập tức, hai bàn tay nhỏ nhắn của con bé vội vã ép chặt lấy hai bên chân váy denim, giữ khít rịt lại để gió lạnh hay bất kỳ cú nẩy nào cũng không thể làm tà váy bay lên được nữa.

Con bé chỉ lẩm bẩm đáp lại: ``\hl{Onii đúng là đồ biến thái\dots{} Quá đáng\dots}'' 

Nhìn cái dáng vẻ vừa giận vừa ngượng, hai tay bám chặt lấy váy bước đi rón rén từng bước của cô em gái, tôi không nén được một tiếng cười khẽ, nhẹ giọng hỏi: ``Anh hỏi thật nhé. Ngoài lý do `trúng tủ nghề' biên đạo với lại lo anh làm việc quá sức ra\dots\ rốt cuộc tại sao em lại muốn nhận việc ở COVER để làm cùng anh đến thế?''

Ryuuhou hơi cúi đầu, đôi mắt dán chặt vào mũi giày sneaker trắng đang bước từng bước chậm rãi. Sau vài giây im lặng, con bé mới khẽ lên tiếng: ``Thật ra\dots\ em nghĩ chuyện này bắt đầu từ lâu rồi.''

Tôi hơi nghiêng đầu nhìn sang, để em tôi tiếp tục.

``Anh còn nhớ hồi trước không? Có một khoảng thời gian anh khó khăn về tiền bạc đến mức phải ngửa tay xin em đấy.''

Tôi gãi đầu, hơi ngượng ngùng. Đúng là hồi ấy tôi từng phải nhờ đến tiền của chính cô em gái mình để xoay xở. Ryuuhou khi đó đã bắt đầu có thu nhập riêng từ hoạt động âm nhạc, còn tôi vẫn chỉ là một cậu con trai ăn bám bố mẹ. Tôi từng nghĩ số tiền ấy có lẽ mình sẽ mất rất lâu mới trả nổi.

``Nhưng anh đã trả rồi còn gì,'' tôi nói, cố giữ giọng bình thản. ``Từ lúc anh bắt đầu nhận lương ở \hll, anh đã trả cho em đàng hoàng, sòng phẳng cả vốn lẫn lãi rồi.''

``Ừ.'' Ryuuhou đáp rất khẽ, và rồi con bé mỉm cười. ``Nhưng anh có biết lúc lần đầu tiên em nhận được tiền anh trả lại, em đã nghĩ gì không?''

Tôi im lặng.

``Em đã không nghĩ là anh sẽ trả được.'' Câu nói ấy khiến bước chân của tôi chậm lại đôi chút.

``K-Không phải vì em nghĩ Onii là người quỵt nợ đâu nhé!'' con bé lập tức chữa lại, hai tay xua xua trước mặt. ``Chỉ là\dots\ lúc đó anh lúc nào cũng tự ép mình phải lo đủ thứ. Tiền ăn, tiền học, tiền sinh hoạt nữa. Em nghĩ số tiền mình đưa anh chắc cứ nằm đó mãi, khi nào anh dư dả hẳn thì trả cũng được.''

Con bé cúi xuống nhìn đôi giày của mình, giọng nhỏ dần: ``Thế mà đến một ngày, anh thật sự đưa nó lại cho em. Đủ từng đồng, thậm chí còn tính cả phần lãi mà anh tự nói là phải trả, mặc dù em cũng không cần anh phải làm vậy.''

Tôi bật cười khẽ. ``Đã vay thì phải trả. Có gì đâu mà em ngạc nhiên.''

``\dots Đúng là như vậy. Với Onii thì có lẽ chẳng có gì đặc biệt.'' Ryuuhou ngẩng đầu lên nhìn tôi. ``Nhưng với em thì khác.''

Con bé khẽ siết hai bàn tay trước ngực. ``Lúc ấy em mới nhận ra là anh thật sự đã trưởng thành rồi. Anh không còn là ông anh cứ phải để em chìa tay ra giúp nữa. Anh đã có thể tự kiếm tiền, tự đứng trên đôi chân của mình, thậm chí còn bắt đầu có thể lo ngược lại cho cả gia đình.''

Tôi im lặng, lắng nghe từng lời em nói. Bỗng nhiên Ryuuhou khựng bước, giọng nhỏ dần như đang sắp xếp lại suy nghĩ trong đầu: ``\dots Onii này, anh có bao giờ thắc mắc vì sao hôm nay em lại chọn bộ đồ này đi họp không?''

Tôi hơi ngước nhìn cô em gái, vừa rồi cô nàng còn đang xấu hổ vì váy bị bay lên, bây giờ lại ngước mắt nhìn thẳng vào tôi, ánh mắt đong đầy sự chân thành khó tả. ``Em vốn chẳng bao giờ thích mặc váy, anh thừa biết mà. Từ nhỏ đến lớn em chỉ thích mặc quần rộng cho thoải mái chạy nhảy. Nhưng hôm nay em cố tình chọn chiếc áo cam tươi, chiếc váy denim này và đôi tất cao cổ\dots\ là vì em muốn bước ra khỏi vùng an toàn của mình. Em muốn xuất hiện trước mặt anh và chị A với một hình ảnh thật sự trưởng thành, chỉn chu, đủ để mọi người tin tưởng giao việc. Em không muốn mãi mãi là đứa em gái nhỏ chỉ biết đứng sau lưng nhận sự bảo bọc của anh nữa.''

Ryuuhou mỉm cười, nụ cười ấy xua tan đi tất cả sự ngượng ngùng ban nãy, đôi mắt sáng rỡ như chứa đầy cả bầu trời: ``Tất cả những gì em làm, tất cả những gì em chuẩn bị từ mấy tháng nay để có thể vào COVER làm việc\dots\ cũng chỉ vì anh thôi, Onii. Em muốn được đứng cạnh anh, san sẻ hết gánh nặng công việc cùng anh. Và em hoàn toàn tự nguyện làm điều đó, không có ai ép buộc em cả.''

Tôi vẫn đứng yên, chưa kịp nói gì thì cô nàng lại tiếp tục, đôi tay khẽ chạm vào cánh tay tôi: ``Em muốn làm việc cùng anh không phải vì em nghĩ Onii yếu đuối đến mức phải có em cứu rỗi. Cũng không phải vì em thương hại anh phải gánh vác quá nhiều. Em chỉ muốn có một ngày nào đó, khi Onii quay đầu nhìn về phía sau, anh sẽ biết rằng lần này, em cũng đang đứng ở đó, sẵn sàng cùng anh vượt qua mọi chuyện.''

Tôi khựng lại, cảm giác ấm áp lan tỏa từ trong tim ra khắp người.

Ryuuhou bật cười nhẹ, kéo nhẹ tay áo tôi để chúng ta tiếp tục bước đi trên vỉa hè: ``Hồi trước anh khó khăn về tiền bạc, em có thể giúp anh bằng tiền. Bây giờ anh cần một người san sẻ công việc, em có thể giúp anh bằng khả năng của em. Đơn giản vậy thôi.''

Con bé nhìn thẳng về phía con phố tấp nập phía trước, nói tiếp: ``Em nhận việc ở COVER chẳng phải vì Onii không tự mình xoay xở được. Em chỉ muốn anh đừng phải tự gánh hết mọi thứ trên vai một mình nữa.''

Tôi nhìn cô em gái nhỏ hơn mình mười tuổi đang bước đều bên cạnh, lòng bất giác thấy ấm lạ. \dots Có lẽ Ryuuhou nói đúng. Ngày trước, tôi từng nghĩ trách nhiệm của một người anh là phải luôn là người cho đi, là điểm tựa vững chắc cho cả gia đình. Nhưng hóa ra, để một gia đình thật sự là một gia đình, đôi khi người anh cũng phải học cách nhận lấy sự giúp đỡ từ những người thương yêu mình ở bên cạnh.

Và có lẽ việc Ryuuhou muốn cùng tôi làm việc cũng chẳng phải để trả lại hay bù đắp cho món nợ năm xưa. Nó chỉ đơn giản là cách cô bé ấy nói với tôi rằng: ``Anh không cần phải làm mọi thứ một mình nữa.''

Tôi đưa tay xoa nhẹ lên mái tóc ngắn của em gái.

``Nào, Onii, đừng có xoa tóc của em mạnh như thế, rối tóc em bây giờ! Trời ạ!'' con bé càm ràm với lấy tay tôi, nhưng rồi cũng bật một nụ cười nhẹ.

Phía trước chúng tôi, con phố Akiba ảo nhộn nhịp trải dài dưới ánh nắng trưa rực rỡ. Tương lai phía trước dĩ nhiên vẫn sẽ còn vô vàn thử thách, nhưng khi nhìn thấy sự trưởng thành từng ngày của các em mình, tôi biết rằng chặng đường sắp tới sẽ không còn là những đêm dài đơn độc nữa, mà sẽ là những ngày tháng đong đầy tiếng cười và sự đồng hành của cả gia đình.

\end{document}